%------------------------------------------------------------------------------%
\documentclass[fleqn,utf8,aspectratio=169,12pt]{beamer}

\usepackage{gaal}

\title{Soma de Vetores}

%------------------------------------------------------------------------------%
\begin{document}

\addtitle

%------------------------------------------------------------------------------%
\section{Soma de Vetores}

%------------------------------------------------------------------------------%
\begin{frame}{Soma de Vetores}
  Dados dois vetores
  \(
    \vec{u}
  \)
  e
  \(
    \vec{v}
  \)
  em \(\R^n\)

  \pause
  Representamos sua soma por
  \[
    \pause
    \vec{w} = \vec{u} + \vec{v}
  \]

  \pause
  A soma representa o ``deslocamento acumulado'' dos dois vetores
\end{frame}

%------------------------------------------------------------------------------%
\begin{frame}{Interpretação Geométrica}
  Para realizar a soma
  \(
    \vec{w} = \vec{u} + \vec{v}
  \)
  geometricamente
  \pause
  \begin{itemize}
      \item colocamos a origem de $\vec{v}$ na extremidade de $\vec{u}$
      \pause
      \item o vetor soma \(\vec{w}\) liga a origem de $\vec{u}$ à extremidade de $\vec{v}$
  \end{itemize}
\end{frame}

%------------------------------------------------------------------------------%
\framedgraphic*{Interpretação Geométrica}{E-03-Soma-1}{}
\framedgraphic*{Interpretação Geométrica}{E-03-Soma-2}{}
\framedgraphic*{Interpretação Geométrica}{E-03-Soma-3}{}

%------------------------------------------------------------------------------%
\begin{frame}{Definição Algébrica da Soma}
  Dados
  \(
    \vec{u} = \vetor{u_1}{u_2}
  \)
  e
  \(
    \vec{v} = \vetor{v_1}{v_2}
  \)

  \pause
  \medskip
  A soma é efetuada componente a componente
  \pause
  \[
    \vec{u}+\vec{v}
    \pause
    \;=\; \vetor{u_1}{u_2} + \vetor{v_1}{v_2}
    \pause
    \;=\; \vetor{u_1+v_1}{u_2+v_2}
  \]
\end{frame}

%------------------------------------------------------------------------------%
%------------------------------------------------------------------------------%
\begin{question}
  Calcule a soma dos vetores
  \[
    \vec{w} = (2, 3)
  \]
  \[
    \vec{v} = (4, -1)
  \]

\end{question}

%------------------------------------------------------------------------------%
\begin{resolution}{Solução}
  \begin{align*}
    \onslide<+->{\vec{w}+\vec{v}}
    \onslide<+->{& = (2, 3) + (4, -1) \\}
    \onslide<+->{& = (2+4,\, 3-1) \\}
    \onslide<+->{& = (6, 2)}
  \end{align*}
\end{resolution}

%------------------------------------------------------------------------------%
\endinput


%------------------------------------------------------------------------------%
\begin{frame}{Propriedades da soma}
  Para quaisquer vetores $\vec{u}$, $\vec{v}$ e $\vec{w}$

  \pause
  \[
    \vec{u}+\vec{v} = \vec{v}+\vec{u}
  \]

  \pause
  \[
    (\vec{u}+\vec{v})+\vec{w} = \vec{u}+(\vec{v}+\vec{w})
  \]

  \pause
  \[
    \vec{u}+\vec{0} = \vec{u}
  \]

  \pause
  \[
    \vec{u}+(-\vec{u}) = \vec{0}
  \]
\end{frame}

%------------------------------------------------------------------------------%
\begin{frame}{Subtração de vetores}
  \onslide<+->{A subtração é definida pela soma com o vetor oposto}
  \begin{align*}
    \onslide<+->{\vec{u}-\vec{v}}
    \onslide<+->{& = \vec{u} + (-\vec{v})}
    \onslide<+->{  = \vec{u} + (-1)\vec{v}                   \\[3mm]}
    \onslide<+->{& = \vetor{u_1}{u_2} + (-1)\vetor{v_1}{v_2} \\[3mm]}
    \onslide<+->{& = \vetor{u_1}{u_2} + \vetor{-v_1}{-v_2}   \\[3mm]}
    \onslide<+->{& = \vetor{u_1-v_1}{u_2-v_2}}
  \end{align*}
\end{frame}

%------------------------------------------------------------------------------%
\begin{frame}{Interpretação geométrica da diferença}
  A diferença
  \(
    \vec{u} - \vec{v}
  \)
  pode ser interpretada como

  \pause
  o vetor que vai da extremidade de $\vec{v}$

  \pause
  à extremidade de $\vec{u}$

  \pause
  quando os dois vetores possuem a mesma origem
\end{frame}

%------------------------------------------------------------------------------%
\framedgraphic*{Interpretação Geométrica}{E-03-Soma-subtracao-1}{}
\framedgraphic*{Interpretação Geométrica}{E-03-Soma-subtracao-2}{}

%------------------------------------------------------------------------------%
\input{ex/E-03-Soma-ex-2}
%------------------------------------------------------------------------------%
\begin{question}
  Dados
  \[
    \vec{u} = \vetor{1}{2}
  \]
  \[
    \vec{v} = \vetor{3}{-1}
  \]
  Determine
  \[
    \vec{w} = 2\vec{u} - 3\vec{v}
  \]
\end{question}

%------------------------------------------------------------------------------%
\begin{resolution}{Solução}
  \[
    2\vec{u}
    \pause
    \;=\; 2\vetor{1}{2}
    \pause
    \;=\; \vetor{2\times 1}{2\times 2}
    \pause
    \;=\; \vetor{2}{4}
  \]

  \pause
  \[
    3\vec{v}
    \pause
    \;=\; 3\vetor{3}{-1}
    \pause
    \;=\; \vetor{3\time 3}{3(-1)}
    \pause
    \;=\; \vetor{9}{-3}
  \]

  \pause
  Portanto
  \[
    \pause
    \vec{w}
    \pause
    \;=\; 2\vec{u} - 3\vec{v}
    \pause
    \;=\; \vetor{2}{4} - \vetor{9}{-3}
    \pause
    \;=\; \vetor{-7}{7}
  \]
\end{resolution}

%------------------------------------------------------------------------------%
\endinput


%------------------------------------------------------------------------------%
%\section{Lista Mínima}
%
%%------------------------------------------------------------------------------%
%\begin{frame}{Lista Mínima}
%  \begin{enumerate}
%    \item
%  \end{enumerate}
%
%  \vfill
%  Atenção: A prova é baseada no livro, não nas apresentações
%\end{frame}

%------------------------------------------------------------------------------%
\end{document}
%------------------------------------------------------------------------------%


